\section*{Source Note: The Hangzhou Generative-AI Decision}
\addcontentsline{toc}{section}{Source Note: The Hangzhou Generative-AI Decision}

Published professional accounts identify the proceeding as \emph{Liang v. Technology Company}, Hangzhou Internet Court, (2025) Zhe 0192 Min Chu No. 18143 (decided Dec. 3, 2025).\lrfootnote{See \bluecite{qiaowenHangzhou2026}.} The comparative discussion relies on Le Weikun's \emph{People's Court Daily} account, hosted on an official Zhejiang government site. A Supreme People's Court Intellectual Property Court webpage reproducing a CCTV interview with the head of the work-report drafting group independently confirms the dismissal, subject-status rule, reasonable-precaution analysis, and absence of proven loss.\lrfootnote{See \bluecite{hangzhouChatbot2026}; \bluecite{spcHangzhouReport2026}.} A complete judgment with a stable, anonymous public URL was not identified. The excerpts below are the author's translation of the court account, not the judgment. They preserve its anonymization, are unofficial, and are supplied for source checking; the Chinese source controls.

\begin{quote}
\small
\textbf{Status and manifestation.} ``AI lacks civil-subject status and cannot make a manifestation of intention.'' The generated compensation promise was not the provider's manifestation because the system could not serve as its messenger, agent, or representative; the provider had not used the model to communicate an intention; ordinary practice did not support reasonable reliance on the random promise; and the provider had not represented that it would be bound by generated content.

\textbf{Classification and governing rule.} The court treated the generative-AI offering as a service under the Interim Measures, not as a product under product-quality law, and applied the Civil Code's general fault rule in article 1165(1).

\textbf{Duty, loss, and causation.} The report describes three layers of care: strict review of legally prohibited information; prominent disclosure of inherent accuracy limits, including immediate warnings where major interests are at stake; and basic reliability precautions using generally accepted technical measures. It reports that the provider displayed limitations in the welcome page, agreement, and interface and used retrieval-augmented generation. The court also found no adequate proof of actual loss or material effect on the user's application decision.
\end{quote}
